\section{Conclusion générale et perspectives}
Le système \textit{Devis IA} propose une chaîne cohérente allant de la demande client au brouillon de devis : compréhension multi-modale, sélection contrôlée des éléments du catalogue, calcul de prix et restitution dans une interface métier. L'usage d'identifiants PostgreSQL réels et la validation des relations entre prestations, options et choix renforcent la fiabilité de la génération.

L'optimisation des instructions avec GEPA fournit un cadre reproductible pour améliorer le module texte. Le choix des modèles ne repose pas sur une impression qualitative : il tient compte de la conformité JSON, de l'identification des prestations, de la quantité, des options, de la latence et de la disponibilité. Les emplacements réservés dans ce rapport permettront d'ajouter les résultats consolidés du benchmark des embeddings, du benchmark final des LLM et des tests applicatifs.

Les perspectives principales sont l'intégration du modèle d'embedding retenu, l'enrichissement du corpus de test réel, la validation terrain des mesures photographiques et la mise en place d'un suivi longitudinal des performances des modèles et des coûts d'inférence.

\newpage
\section*{Références à compléter}
\begin{enumerate}
\item Documentation GEPA et protocole d'optimisation utilisé dans les expérimentations.
\item Publication ou dépôt de référence ExtractBench concernant l'appariement hongrois des structures répétées.
\item Source des benchmarks publics utilisés pour le choix du modèle ASR français Hojo.
\item Documentation des fournisseurs de modèles et quotas utilisés dans le benchmark final.
\end{enumerate}
